% Created 2022-04-08 Fri 22:21
% Intended LaTeX compiler: pdflatex
\documentclass[a4paper, 11pt, twoside]{article}
\usepackage[utf8]{inputenc}
\usepackage[T1]{fontenc}
\usepackage{graphicx}
\usepackage{grffile}
\usepackage{longtable}
\usepackage{wrapfig}
\usepackage{rotating}
\usepackage[normalem]{ulem}
\usepackage{amsmath}
\usepackage{textcomp}
\usepackage{amssymb}
\usepackage{capt-of}
\usepackage{hyperref}
\usepackage{baskervillef}
\usepackage{geometry}\geometry{ a4paper, total={170mm,257mm}, left=20mm, top=20mm,}
\usepackage{hyperref}\hypersetup{pdfauthor={Olivier Rey}, pdftitle={Ressources pour JDR}, pdfkeywords={jdr, ressources, orey-jdr}, pdfsubject={jeu de rôles}, pdfcreator={Emacs 26.1 (Org mode 9.1.9)}, pdflang={French}, colorlinks=true, linkcolor={blue}, urlcolor={blue}}
\usepackage{titlesec}\titlelabel{\thetitle. \quad}
\usepackage[table,svgnames]{xcolor}\rowcolors{1}{Gainsboro}{WhiteSmoke}
\usepackage{etoolbox}\AtBeginEnvironment{longtable}{\small}
\author{rey.olivier@gmail.com}
\date{2022-04-08}
\title{Ressources pour JDR}
\hypersetup{
 pdfauthor={rey.olivier@gmail.com},
 pdftitle={Ressources pour JDR},
 pdfkeywords={},
 pdfsubject={},
 pdfcreator={Emacs 26.1 (Org mode 9.1.9)}, 
 pdflang={Frenchb}}
\begin{document}

\maketitle
\tableofcontents

\begin{center}
\includegraphics[width=4cm]{logo-orey.png}
\end{center}

Les différents dossiers de ce repo contiennent des ressources en français pour JDR ou des traductions (voir aussi la \href{RessourcesPourJDR-ORey.pdf}{version PDF de cette page}).

Ce repo est le premier d'un cycle de 3. Voici les deux suivants :
\begin{itemize}
\item \href{https://github.com/orey/DandD}{github.com/orey/DandD} : Un repo en anglais dédié à D\&D, tout spécialement aux vieilles choses (0e et 1e),
\item \href{https://github.com/orey/ttrpg}{github.com/orey/ttrpg} : Un repo en anglais dédié à de nombreux jeux de rôles et contenant une grande liste de liens.
\end{itemize}

\section{Dernières livraisons}
\label{sec:org5cf657d}
\begin{longtable}{cp{2cm}p{1.5cm}p{4cm}ccp{4cm}}
\textbf{Date} & \textbf{JDR-RPG} & \textbf{Type} & \textbf{Livrable} & \textbf{Format} & \textbf{Ref} & \textbf{Commentaire}\\
2022 & \href{https://github.com/orey/jdr/tree/master/Mythic-fr}{Mythic} & GM Emulator & \href{https://github.com/orey/jdr/blob/master/Mythic-fr/MythicGME-EcranMJ-VersionFrancaise-OreyJdr05.pdf}{Ecran en français pour Mythic GME} & PDF & 05 & Traduction originale\\
2021 & \href{https://github.com/orey/jdr/tree/master/Fudge-fr}{Fudge} Department 13 & JDR & \href{https://github.com/orey/jdr/blob/master/Fudge-fr/Division13/Department13-FeuillePerso.pdf}{Feuille de perso Department 13} & PDF &  & Pour le setting Department 13\\
2021 & \href{https://github.com/orey/jdr/tree/master/Fudge-fr}{Fudge} & JDR & \href{https://github.com/orey/jdr/blob/master/Fudge-fr/FudgeEnUnePage/FudgeEnUnePage-OReyJdr04.pdf}{Fudge en une page} & PDF & 04 & Traduction originale\\
2021 & \href{https://github.com/orey/jdr/tree/master/FightingFantasys-fr}{Fighting Fantasy VF} & JDR & \href{https://github.com/orey/jdr/blob/master/FightingFantasys-fr/FightingFantasy-VersionFrancaise-OreyJdr03.pdf}{Fighting Fantasy Version Française} & PDF & 03 & Traduction et adaptation originale\\
2021 & Tous & Aide de jeu & \href{https://github.com/orey/jdr/blob/master/G\%C3\%A9n\%C3\%A9rateurLabyrinthe/GenerateurDeLabyrinthe-OreyJdr02.pdf}{Générateur de labyrinthe pour JDR} & PDF & 02 & Traduction et adaptation originale\\
2021 & \href{https://github.com/orey/jdr/tree/master/DungeonSquad-fr}{DungeonSquad! VF} & JDR pour enfants & \href{https://github.com/orey/jdr/blob/master/DungeonSquad-fr/DungeonSquad-VersionFrancaise-OreyJdr01.pdf}{DungeonSquad! Version Française} & PDF & 01 & Traduction et adaptation originale\\
2021 & \href{https://github.com/orey/jdr/tree/master/DungeonSquad-fr}{DungeonSquad! VF} & JDR pour enfants & \href{https://github.com/orey/jdr/blob/master/DungeonSquad-fr/DungeonSquadFr-FeuillePerso.pdf}{Feuille de perso} & PDF &  & Pour fille et garçon\\
2021 & \href{https://github.com/orey/jdr/tree/master/DungeonSquad-fr}{DungeonSquad! VF} & JDR pour enfants & \href{https://github.com/orey/jdr/blob/master/DungeonSquad-fr/DungeonSquadFr-Ecran.pdf}{Ecran} & PDF &  & Un outil indispensable\\
2021 & \href{https://github.com/orey/jdr/tree/master/AppelDeCthulhu}{Appel de Cthulhu} & Horreur & \href{https://github.com/orey/jdr/blob/master/AppelDeCthulhu/AppelDeCthulhu-EcranComplementaire.pdf}{Ecran complémentaire MJ} & PDF &  & Ecran complémentaire MJ\\
2021 & \href{https://github.com/orey/jdr/tree/master/AppelDeCthulhu}{Appel de Cthulhu} & Horreur & \href{https://github.com/orey/jdr/blob/master/AppelDeCthulhu/AppelDeCthulhu-FicheDeTemps.pdf}{Fiche de temps} & PDF &  & Pour l'Appel de Cthulhu ou autre jeu Basic RPS\\
2021 & Tous & Aide de jeu & \href{https://github.com/orey/jdr/blob/master/Aftermath/LocalisationDesBlessures.png}{Localisation des blessures} & PNG &  & A intégrer dans une synthèse d'aides de jeu pour MJ\\
2021 & \href{https://github.com/orey/jdr/tree/master/Risus-fr/}{Risus} & Flowchart & \href{https://github.com/orey/jdr/blob/master/Risus-fr/risus-flowchart.pdf}{Flowchart complet du jeu} & PDF &  & Peut servir d'éran\\
2021 & \href{https://github.com/orey/DandD}{D\&D 5e} & Couverture & \href{https://github.com/orey/DandD/blob/master/DandD\_5e\_BasicEditionLuluCover/Cover.pdf}{Couverture pour D\&D 5e Basic Rules} & PDF &  & Pour Lulu.com\\
\end{longtable}

\section{Liens}
\label{sec:orgea00e2a}

\subsection{Sites de jeux en français}
\label{sec:org3751ff8}

\begin{longtable}{p{7cm}p{7cm}}
\textbf{Type} & \textbf{Site}\\
\textbf{C} & \\
Le cénotaphe & \url{http://casquenoir.free.fr/index.php}\\
Créatures légendaires & \url{https://fr.wikipedia.org/wiki/Liste\_de\_cr\%C3\%A9atures\_l\%C3\%A9gendaires}\\
\textbf{D} & \\
Blog de Jérôme Darmont & \url{http://darmont.free.fr/}\\
Discussions de Rôlistes Ouvertes et Libres & \url{https://www.facebook.com/groups/254213402190606}\\
\textbf{E} & \\
Echecs: Check \& Strategy, site en français & \url{https://www.chess-and-strategy.com}\\
Empire Galactique JDR, un classique & \url{https://sites.google.com/site/empiregalact}\\
Egrégore, un JDR fantastique & \url{https://business.facebook.com/EgregoreJdR/?business\_id=456290144533916}\\
Epées et Sorcellerie JDR & \url{https://sites.google.com/site/wizardinabottle/epeesetsorcellerie}\\
\textbf{F} & \\
Une traduction française du RPG "FU" & \url{https://brunobord.gitbooks.io/fu-rpg-libre-et-universel/}\\
Faenix & \url{https://chezfaenyx.blogspot.com/2021/11/20-jeux-20-questions-episode-3.html}\\
Traduction française de Fudge & \url{http://fudge.ouvaton.org/}\\
\textbf{G} & \\
Giannirateur de scénarios & \url{http://loukoum.online.fr/jdr/adj/gianni1.htm}\\
 & \url{http://loukoum.online.fr/jdr/scenars/defi2012.htm\#47}\\
Le Grog, répertoire de JDR et news & \url{http://www.legrog.org/}\\
\textbf{H} & \\
Harry Potter JDR, un très beau travail & \url{https://www.geek-it.org/harry-potter-jdr}\\
Heroquest, un site de fan & \url{https://www.heroquest-revival.com}\\
\textbf{I} & \\
Imaginos & \url{https://blogs.bl0rg.net/imaginos/}\\
L'insoutenable légèreté du joueur & \href{http://linsoutenablelegeretedujoueur.over-blog.com/}{linsoutenablelegeretedujoueur.over-blog.com}\\
\textbf{K} & \\
Koma JDR et autres jeux de Xavier Raoult & \url{http://komajdr.free.fr/?page\_id=96}\\
\textbf{L} & \\
Les jeux d'Olivier Legrand & \url{http://storygame.free.fr/}\\
\textbf{M} & \\
Maléfices vieux suppléments & \url{https://www.scribd.com/user/381722775/Jean-Charles-BLANGENOIS}\\
Méga 4, un site avec des suppléments exceptionnels ! & \href{https://www.messagers-galactiques.com}{messagers-galactiques.com}\\
\textbf{O} & \\
Osric JDR & \url{https://osric.fr}\\
\textbf{P} & \\
Les meilleurs articles TTRPG traduits & \href{https://ptgptb.fr}{ptgptb.fr}\\
\textbf{V} & \\
La voix d'Héort, ressources pour Glorantha & \url{https://heort.wordpress.com/}\\
 & \\
\end{longtable}

\subsection{Magazines en français}
\label{sec:org55d5dcc}

\begin{longtable}{p{7cm}p{7cm}}
\textbf{Type} & \textbf{Site}\\
\textbf{B} & \\
Les anciens "Backstab" & \url{https://www.abandonware-magazines.org/affiche\_mag.php?mag=199}\\
\textbf{C} & \\
Les anciens "Casus Belli" & \url{https://www.abandonware-magazines.org/affiche\_mag.php?mag=188}\\
\textbf{G} & \\
Quelques vieux "Graal" & \url{https://www.abandonware-magazines.org/affiche\_mag.php?mag=402}\\
\textbf{J} & \\
Les anciens "Jeux et Stratégie", un must & \url{https://www.abandonware-magazines.org/affiche\_mag.php?mag=185}\\
\textbf{T} & \\
Les vieux "Tangente" & \url{https://www.abandonware-magazines.org/affiche\_mag.php?mag=326}\\
 & \\
\end{longtable}

\section{Explorations récentes}
\label{sec:org42de0d3}

\begin{longtable}{cp{2cm}p{1.5cm}p{7cm}ccc}
\textbf{Date} & \textbf{Game} & \textbf{Type} & \textbf{Comment} & \textbf{Note} & \textbf{OSR} & \textbf{Ongoing}\\
2022 & Méga & Sci-Fi & Un jeu fantastique par son univers. Procurez-vous l'univers de Méga 4 ! & \textbf{5/5} & N & \textbf{Y}\\
2022 & Risus & Generic system & Bon, je me ravise, c'est un bon jeu. Se procurer l'écran pour les probas. & 4/5 & N & \textbf{Y}\\
2022 & AD\&D 1e, version US & Heroic Fantasy & Un travail colossal, un Gary Gygax au sommet de sa forme ! La matrice du JDR ! & \textbf{5/5} & - & \textbf{Y}\\
2022 & \href{https://www.cortexrpg.com/compendium/explore-the-rules/}{Cortex} & Generic System &  & - & N & Later\\
2022 & \href{https://www.drivethrurpg.com/product/117563}{The Void} & Horror Sci-Fi & Interesting Cthulu Saga in space & - & N & N\\
2022 & Troika! & Generic system & A reinterpretation of the \href{https://github.com/orey/jdr/tree/master/FightingFantasys-fr}{Fighting Fantasy} rules with funny elements & 4/5 & Y & \textbf{Y}\\
2022 & 1PG Star Legion & Sci-Fi & A sci-Fi small RPG & - & N & N\\
2022 & \href{https://www.drivethrurpg.com/product/186894/Cepheus-Engine-System-Reference-Document}{Cepheus engine} & Sci-Fi & The SRD of the Traveller TTRPG & - & N & N\\
2022 & \href{https://www.chaosium.com/runequest-starter-set/}{Runequest Starter Set} & Heroic Fantasy & A great game & \textbf{5/5} & N & N\\
2022 & Hurlements (1989) & Middle-Age & Strange French game, at the center of the narrativist French trend & 2/5 & N & N\\
2021 & \href{https://github.com/orey/jdr/tree/master/BladesInTheDark-SRD}{Blades in the Dark SRD} & Heroic Fantasy &  & - & N & Later\\
2021 & \href{https://github.com/orey/jdr/tree/master/Risus-fr}{Risus} & Generic system & Irony with Clichés and D6 with difficulty factors & 3/5 & N & N\\
2021 & \href{https://www.drivethrurpg.com/product/89534/FU-The-Freeform-Universal-RPG-Classic-rules}{FU} & Generic system & Very basic system for roleplay & 3/5 & N & N\\
2021 & \href{http://www.onesevendesign.com/laserfeelings/}{Lasers and Feelings} & Sci-Fi & Great simple RPG & 4/5 & N & N\\
2021 & GURPS & Generic system & A great classical system with great supplements & 4/5 & N & Later\\
2021 & \href{https://github.com/orey/jdr/blob/master/Fudge-fr/FudgeEnUnePage-ORey03.pdf}{Fudge} (en une page) & Generic system & An "open GURPS" with a 7-levels ladder and scales. Very GURPS inspired & 4/5 & N & Later\\
2021 & \href{http://komajdr.free.fr/fichiers/BiTs.rar}{Bits } & Generic system & In French, a one page generic system & - & N & N\\
2021 & \href{http://storygame.free.fr/}{Trucs trop bizarres} & Modern kids & In French, a very simple game system & 3/5 & N & N\\
2021 & Advanced Fighting Fantasy & Heroic Fantasy & To play with children & - & N & Later\\
2021 & Modern AGE system & Modern & Ongoing & - & N & Later\\
2021 & Tunnels \& Trolls 1e & Heroic Fantasy & Interesting & 4/5 & N & N\\
2021 & Alternity 98 & Modern (Generic) & A very good system abandonned by WotC for the crappy D20 Modern & \textbf{5/5} & N & Later\\
2021 & The Esoterrorists 2e & Modern & The first Gumshoe system & - & N & Later\\
2021 & The Dragon & Press & Old issues of The Dragon, in \href{https://archive.org/details/DragonMagazine045\_201903}{archive.org} (1-100 251-280) & - & - & N\\
2021 & D20 Modern SRD & Generic System & Exploration in parallel to some \href{https://archive.org/details/Polyhedron105}{Polyhedron} readings & 2/5 & N & N\\
2021 & Gumshoe system SRD & Generic System & Entering into simplified translation process & - & N & Later\\
2021 & 13th Age & Heroic Fantasy & Just starting & - & Y & Later\\
2021 & Basic Roleplaying System & Generic System & The best, especially for CoC, free ed. is great & \textbf{5/5} & N & Later\\
2021 & The Wretched & Horror & Bof & 2/5 & N & N\\
2021 & GURPS & Generic System & Not convinced & 4/5 & N & N\\
2021 & Fighting Fantasy & Generic System & From Steve Jackson \& Ian Livingstone : \href{https://github.com/orey/jdr/tree/master/FightingFantasys-fr}{French translation} & 4/5 & Y & N\\
2021 & Bloodlust & Heroic Fantasy & French game by Croc & 3/5 & N & N\\
2021 & Metamorphosis Alpha & Sci-Fi & Interesting game & 3/5 & - & N\\
2021 & Ironsworn & Heroic Fantasy & Interesting game but too random (action dice vs 2D10) & 3/5 & N & N\\
2021 & Gumshoe system & Generic system & Investigation oriented: That one is for me :) & - & N & Later\\
2021 & DCC & Heroic Fantasy & A whole universe & 4/5 & Y & N\\
2021 & Légendes & Historic Fantasy & Great game for the universes. Hyper complex game system & 4/5 & N & Later\\
2021 & Tékumel & Heroic Fantasy & Author's world & 3/5 & N & N\\
2021 & Microlite & Generic System & \href{https://github.com/orey/jdr/tree/master/Microlite20-fr}{French translation} done. Not playable as-is. & 3/5 & N & N\\
2021 & \href{http://www.fortuneswheel.co.uk/}{Fortunes Wheel} & Witching Tales & Very interesting with tarot cards & 4/5 & N & Later\\
2021 & Maléfices & French Steampunk & Un des meilleurs JDR français & \textbf{5/5} & N & Later\\
2021 & GURPS & Generic System & To investigate & - & N & N\\
2021 & Traveller 1e & Sci-Fi & Seducing & - & N & Later\\
2020 & D\&D 5e basic rules & Heroic Fantasy &  & 3/5 & - & N\\
2020 & Covetous & GM Emulator & Bon produit avec plein de tables & - & N & Later\\
2020 & Conspiracy X & Modern &  & - & N & Later\\
2020 & D\&D SRD 3.5 & Heroic Fantasy & \href{https://github.com/orey/srd-3.5}{Repo spécial} avec diverses versions. & 4/5 & - & N\\
2020 & Méga & Sci-Fi & A French success & - & N & Later\\
2020 & Empire galactique & Sci-Fi & One of the first french RPG & 3/5 & N & N\\
2020 & L'appel de Cthulhu & Horror & The best & \textbf{5/5} & N & Later\\
2020 & Warhammer FR 1e & Heroic Fantasy & A very good game, surtout pour la Campagne Impériale & \textbf{5/5} & N & Later\\
2020 & Hero kids & RPG for kids & Bof, better play a simple adult game, or Bubblegumshoe & 2/5 & N & N\\
2020 & Pokethulhu & Fun & You need to like the comics & 2/5 & N & N\\
2020 & CRGE & GM Emulator & Based on the "Yes but\ldots{}/No but\ldots{}" & 2/5 & N & N\\
2020 & Mythic & GM Emulator & Great! \href{https://github.com/orey/jdr/tree/master/Mythic-fr}{Resources in French} (un écran !) & \textbf{5/5} & N & Later\\
2020 & PIP system & Generic system &  & - & N & Later\\
2020 & QAGS - Quick Ass Game System & Generic system & Simple and funny dynamic system & 4/5 & N & Later\\
2020 & Gateway & Heroic fantasy & Based on D\&D & 2/5 & Y & N\\
2020 & FU - Freeform Universal & Generic system & JDR basé sur le "Yes but\ldots{}/No but\ldots{}" & 3/5 & N & N\\
2020 & \href{https://github.com/orey/jdr/tree/master/Risus-fr}{Risus} & Generic system & In French:  \href{https://github.com/orey/jdr/tree/master/Risus-fr}{Règles résumées Risus} avec flowchart & 3/5 & N & N\\
2020 & PremièreFable (FirstFable) & JDR pour enfants & Traduction de FirstFable. Lien : \href{https://orey.github.io/premierefable/}{PremièreFable le JDR}. & 4/5 & N & N\\
2020 & \href{https://www.drivethrurpg.com/product/144558/Mini-Six-Bare-Bones-Edition}{MiniSix} & Generic system & D6 & - & N & Later\\
2020 & Dagger & RPG for kids & Bof & 2/5 & Y & N\\
\end{longtable}

\section{Quelques réflexions sur les systèmes de jeux}
\label{sec:orgb4ce5ec}

Ci-dessous, quelques réflexions les systèmes de jeux et autres marronniers du JDR.

\subsection{Powered by the Apocalypse (PbtA)}
\label{sec:orgc0793d1}

Certains jeux récents, notamment la vague PbtA (\href{https://en.wikipedia.org/wiki/Powered\_by\_the\_Apocalypse}{Powered by the Apocalypse}), reprennent à leur compte des questions vieilles comme le JDR (du roleplay ou des règles, disait-on dans le temps) pour leur apporter des "solutions" pour le moins particulières. Ces solutions méritent un examen détaillé.

\subsubsection{Historique}
\label{sec:org1a56790}

La vague PbtA a commencé par des discussions sur le forum \href{https://github.com/orey/ttrpg\#the-forge}{The Forge}. En quelque sorte, PbtA est la troisième itération d'un \emph{modèle théorique du jeu de rôles}, les deux premières étant :
\begin{itemize}
\item En 1997, le \href{https://en.wikipedia.org/wiki/Threefold\_model}{Threefold model} (Dramatist, Gamist, Simulationist), apparu sur un forum de discussions,
\item Entre 1999 et 2005, la \href{https://en.wikipedia.org/wiki/GNS\_theory}{GNS Theory} (Gamism, Narrativism, Simulationism), portée par un gourou, \href{https://en.wikipedia.org/wiki/Ron\_Edwards\_(game\_designer)}{Ron Edwards}.
\end{itemize}

Ces modèles conceptuels du jeu de rôle, relativement étroits car tri-dimensionnels, ont nourri une réflexion permettant la construction d'autres types de jeux. La vague PbtA est, aujourd'hui, le grand vainqueur de cette histoire.

\subsubsection{Caractéristiques critiquées sur les JDR "anciens"}
\label{sec:orga566578}

Voici quelques uns des arguments que l'on trouve sur les forums concernant le pourquoi de la vague PbtA et les limitations des jeux anciens.

\textbf{Complexité des règles}

Les règles sont trop complexes, trop simulationnistes, trop spécifiques (une règle différente par cas sans unité globale), trop incohérentes (pas de ligne directrice globale), trop "crunchy" comme disent les américains. Les règles des anciens jeux font trop appel à des lancers de dés incessants, et à des modificateurs qui s'empilent de manière complexe, à des centaines de pages de règles.

\textbf{Directivité du MJ}

Le MJ est trop directif et il ne met pas en place un univers collaboratif où les joueurs peuvent co-construire l'univers avec lui au travers de l'utilisation de la discussion.

\textbf{Préparation du jeu}

Le temps de préparation est trop long, trop important, pour les anciens jeux. L'investissement du MJ est trop important. L'investissement demandé n'est plus adapté à notre monde moderne, que ce soit pour le MJ ou pour les joueurs.

Derrière cette préoccupation matérielle, on retrouve la co-construction de l'histoire et une volonté de créer une \emph{dynamique organique} du jeu. Le jeu se développerait "de lui-même" et n'aurait donc pas besoin de guidelines voire de scénarios menés par un MJ ayant tout péparé seul.

\textbf{Non intégraion des dernières techniques de jeu}

Certaines variantes de jeu ont permis de voir les choses différemment (JDR solo avec "gamemaster emulator", JDR sans MJ, etc.) et de pousser les jeux "narratifs" et co-construits sur le devant de la scène. Les anciens jeux n'ont pas inclus ces nouvelles tendances.

\textbf{Des arguments anciens}

Les arguments critiques envers les jeux très lourds en règles (à commencer par D\&D) étaient à peu près les mêmes, durant les années 80/90. La recherche du Graal du jeu basé sur le \emph{roleplay} et avec très peu de règles a longtemps occupé les colonnes des magazines de JDR. de nombreuses tentatives ont été faites, sans être vraiment couronnées de succès.

\textbf{Un nouveau type de jeu}

Apocalypse World, le premier jeu de la vague PbtA sort de 2010, soit plus de 20 ans après ces débats. La famille Baker, qui a designé ce jeu, est une famille de quadras à l'époque. Nourris par The Forge et la GNS, ils proposent un nouveau type de jeu.

\subsubsection{Une modification profonde de la vision des PJ}
\label{sec:orgd410319}

\begin{enumerate}
\item Des PJ à base de chiffres
\label{sec:org63c9847}

La première composante du JDR traditionnel est la composante PJ. Selon comment ces derniers sont structurés, les joueurs auront plus ou moins de possibilités. Les règles du jeu organisent les interactions entre les PJ et le monde, ainsi qu'avec les PNJ.

Les jeux de rôles de la première génération (D\&D, Cthulhu, GURPS, Rolemaster, etc.) sont basés sur une possibilité de \emph{comparaison objective} entre les PJ, et sur un calibrage des PJ par rapport à la réalité du monde. Ainsi, dans la plupart des jeux, on trouve, décrit en termes de jeu, les principaux animaux. Etablir une échelle objective permet aussi de construire et de dimensionner des adversaires PNJ. Certains jeux comme D\&D utiliseront la notion de niveau (voir quelques \href{https://github.com/orey/DandD}{commentaires sur D\&D}).

\item Des PJ plus ou moins structurés avec des mots ou des expressions
\label{sec:org251274d}

Avec l'apparition de jeux comme \href{https://github.com/orey/jdr/tree/master/Fudge-fr}{Fudge} (1992), un "pont" est dressé entre des valeurs sous forme de nombre (de -3 à +3) et des \emph{descriptifs qualificatifs} portant cette "valeur". Même si la mécanique sous-jacente est encore à base de nombres et de modificateurs, les mots vont prendre progressivement une importance énorme dans le monde des JDR, jusqu'à prendre la place d'attributs, de compétences, de dons, de défauts ou même de pouvoirs.

Les mots s'imposent dans les créations des PJ au travers des "aspects", des "clichés", des "archétypes", des "avatars", des "concepts de personnages", etc. Ces mots peuvent être invoqués durant le jeu pour appeler une mécanique particulière, le plus souvent un bonus ou une compétence impliquée par ou "sous-entendue" dans l'expression elle-même.

Or, au travers de cette irruption des mots dans les mécaniques des JDR, les problèmes relatifs à l'ambiguïté des mots et des expressions entrent dans le monde du JDR.

\textbf{Des jeux pour vétérans ?}

Il faut une certaine expérience du JDR pour pouvoir définir des \emph{mots et expressions utiles} au jeu. En un sens, le JDR s'adresse implicitement à des vétérans, voire à des vétérans dans un mode jeu ironique (voire cynique). Ces derniers ont les codes du JDR : ils sauront quelle expression leur permet de "joueur comme à AD\&D".

Le discours est que les jeux utilisant des mots rendent le JDR plus accessible aux débutants. Etant donné que la logique du jeu est plus complexe à comprendre, car souvent implicite, cette assertion n'est pas facilement défendable.

\textbf{Incertitude autour des PJ et des PNJ}

Les mots ou phrases introduisent une \emph{incertitude autour des personnages}, incertitude à laquelle le MJ doit s'adapter.

En effet, les mots sont vagues, soumis à des interprétations et parfois en recouvrement sémantique partiel, ce qui rend compliqué leur usage. Cette ambiguïté est vue comme positive par les tenants de cette mécanique. Elle bénéficie clairement aux joueurs en dépossédant le MJ d'une partie de ses responsabilités.

Les mots et les phrases descriptives des PJ et PNJ  introduisent aussi une \emph{incertitude dans l'équilibre des PJ entre eux et avec les PNJ}, et donc une possibilité d'arbitraire pour les MJ. Il devient, en effet, complexe de quantifier les expressions pour équilibrer les personnages et les aventures.

\textbf{Favoriser les extravertis}

Les mots et les phrases \emph{favorisent les joueurs extravertis} qui pourront interpréter de manière libre un "cliché" alors que les introvertis seront desservis par des règles basées sur des mots et sur le besoin d'improviser oralement pour les interpréter. Dans les jeux classiques, les joueurs utilisent des échelles objectives, ce qui permet à tous d'avoir une échelle commune pour s'exprimer.

\textbf{Deux types de MJ}

Les jeux à base de mots sont des jeux très spécifiques dans lesquels le MJ doit investir pour comprendre comment utiliser ces mots ou expressions plus ou moins libres et plus ou moins contraintes par le game designer. Cet investissement est souvent lié à la recherche d'un mécanisme ressemblant à une échelle objective.

D'autres MJ, en revanche, se satisferont du côté plus "libre" d'un jeu basé sur les mots et sur les interprétations des joueurs et du MJ. Peut-être est-on face à deux types de MJ : des MJ scientifiques aimant les échelles absolues et les MJ littéraires aimant l'ambiguïté des mots.

\item La structuration du "backgound" des personnages
\label{sec:org5700836}

Le JDR narratif apporte une obsession étrange : celle de la \emph{structuration} du background. Dans la plupart des jeux anciens, même s'il était recommandé de créer un background à son personnage, cela était plus ou moins fait, et disons que le background s'enrichissait au fur et à mesure des parties. D'une certaine façon, même dans les jeux modernes, les aventures passées des PJ devenaient leur backgound principal.

Beaucoup de jeux plus récents établissent une vraie \emph{dictature du background} en exigeant de le structurer de manière schématique voire caricaturale, en termes de jeu. Ainsi, on voit apparaître, en plus de la notion de "cliché" ou d'"archétype"  :
\begin{itemize}
\item Une certaine obsession pour les \emph{défauts des PJ}, souvent utiles pour gagner des points dans le processus de création, parfois obligatoires dans le processus, souvent en contrepartie des dons ;
\item Des contacts sociaux obligatoires,
\item Un ennemi juré obligatoire,
\item Etc.
\end{itemize}

Cette schématisation à outrance du background des personnages concourt à en faire des caricatures, semblant être issues du même moule, et à rendre suffocant l'univers des personnages. Dans une conception ancienne, le PJ doit être libre avant tout et n'a pas à être "backgroundisé" arbitrairement au travers de dimensions caricaturales.

Ce sujet est délicat, car il s'agit d'un problème de curseur. La structuration du background des personnages a toujours été un sujet dans le JDR, mais dans des jeux comme PbtA, il semble que le bouchon soit poussé plus loin.

\item Les playbooks : des caricatures de classes de PJ ?
\label{sec:org95d0f6a}

La notion de classe ne disparaît pas dans les jeux de type PbtA, on pourrait même dire qu'elle s'hyper-spécialise. Au travers des "playbooks", il est difficile de savoir si l'on joue un personnage prétiré ou un personnage d'une classe de personnage très "étroite". D'autant qu'au sein d'une partie, un seul personnage d'un certain type ne peut être joué dans le groupe.

Le besoin de customiser le playbook est donc limité à sa plus simple expression, car le personnage est déjà si spécifique, qu'il ne mérite qu'un "acteur" pour l'incarner. D'une certaine façon, la logique des playbooks équivaut, pour les jeux plus anciens, à une logique de personnages prétirés. Garantir la bonne répartition des rôles implique, tout comme dans D\&D, équilibrer un groupe pour une aventure.

\item Des caricatures de PJ pour un petit univers ?
\label{sec:org43424a2}

Si l'on mèle description textuelle et playbooks, nous voyons un JDR où les personnages ont tendance à être des caricatures.

La volonté de jouer des caricatures semble clairement influencée par :
\begin{itemize}
\item Les films et séries,
\item Les jeux vidéos,
\item Les comics et les mangas.
\end{itemize}

Or, jouer des caricatures ou des archétypes de personnages de films ou de BD offre, pour certains joueurs et MJ, des possibilités très limitées. Ou disons que cette approche permet de jouer des aventures ressemblant à des films ou des séries, mais sur des styles souvent très étroits, voire très très étroits.

Cette \emph{réduction de l'univers de jeu} semble une caractéristique des jeux PbtA. Les PJ archétypaux évoluent dans un monde restreint, archétypal. Notons que cette caractéristique implique la multiplication des jeux, nous y reviendrons.

Dans des jeux plus traditionnels, les règles permettent de gérer l'univers de jeu qui, par nature, est vaste, voire - on aime à le penser - infini. La partie "liberté des personnages à interagir avec le reste du monde" est une dimension importante du jeu, d'où les cartes des mondes avec des détails à différentes échelles. Le MJ est garant de ce monde qui, en aucun cas, ne serait se limité aux quelques archétypes du genre.

Le style même des aventures peut changer d'une aventure à l'autre. Si une aventure est axée sur l'exploration et les combats, une seconde peut être une enquête et une troisième peut avoir un scénario de film d'horreur.

En fait, c'est comme si les jeux de super-héros avaient gagné en esprit sur le reste du JDR : il semble qu'il faille jouer des caricatures standardisées dans les jeux modernes, des personnages de série télé, qui ont été volontairement fabriqués pour évoluer dans un certain genre de séries très étroit (sous-genre ou sous-sous-genre).
\end{enumerate}

\subsubsection{Les "move"}
\label{sec:org1ac360c}

Dans certains jeux relativement récents, les mots sont utilisés pour les "move", des genres d'actions génériques que les PJ peuvent faire, ces actions étant adaptées au sous-genre proposé par le jeu.

Le game designer impose, en plus de tout ce que nous avons vu, une certaine manière de jouer les personnages en déterminant les "move" par type de personnage, pour tous les types, ainsi que pour le MJ. Généralement, ces moves sont spécifiques au jeu, peu faciles à comprendre et à utiliser, et en recouvrement sémantique les uns avec les autres (ce qui est parfois volontaire dans certains jeux).

Si l'on examine finement des jeux comme \emph{Apocalypse World}, on se rend compte que les moves des PJ et du MC (Master of Ceremony, un genre de "MJ light et collaboratif") semblent se répondre et avoir fait l'objet de simulations par ordinateur. En contraignant les actions possibles des PJ, et en catégorisant les différentes actions possibles du MC, le jeu donne l'impression de mettre en place une \emph{chorégraphie} bien adaptée à certains jeux ou univers, mais au final très mécaniste, et dont les résultats ont été "optimisés" dans des simulateurs inaccessibles aux joueurs et au MJ.

Certes, le jeu possède cette \emph{dynamique organique} dont nous parlions auparavant, et donc fournit les effets promis. Mais cela est rendu possible au travers des enchaînements de moves entre PJ et MC, ces dernières ayant fait l'objet de simulations informatisées ayant été utilisées dans le game design.

Les moves ne sont donc pas des vraies compétences, mais plus, comme leur nom l'indique, des actions appelant d'autres moves. Le glissement sémantique est subtil, mais réel : dans une situation, l'utilisation d'un move appelle un autre move "compatible" en face, et donc restreint mécaniquement les possibilités de jeu, en favorisant des "modèles d'action" (\emph{action patterns}) archétypaux au genre considéré.

\subsubsection{Théâtre et jeu de plateau}
\label{sec:org1d9d73a}

Nous sommes donc devant une nouvelle catégorie de jeux ayant les caractéristiques suivantes.

\textbf{Les personnages sont des \emph{archétypes} (playbook) aux actions archétypales (move)}

Pour les anciens rôlistes, ce sont des "personnages prétirés". Les playbooks vont au delà des classes de personnages de D\&D pour qui la classe est un moyen de progresser dans une certaine voie.

Les aventures sont étudiées et optimisées pour ces archétypes ayant ces actions archétypales. Le joueur doit endosser l'archétype, jouer le personnage prétiré et utiliser ses actions (move) prédéfinis. S'il ne le joue pas, l'histoire et sa dynamique peuvent être perturbées.

Nous sommes donc dans une perspective plus \emph{théâtrale} que libre. Le canevas imposé aux joueurs semble plus dur que le canevas imposé par les anciens jeux.

\textbf{Le JDR est devenu un jeu de plateau}

Si les moves devenaient des cartes physiques et les playbooks des fiches cartonnées, nous pourrions être dans un jeu de plateau, type jeu de stratégie. Une scène devient alors une série de cartes à action que chaque jouer abat à son tour (PJ et MC ou PJ et PJ).

Les jeux de type PbtA ne sont donc pas des JDR classiques mais ressemblent plus à des jeux hybrides entre jeu de rôle, jeu de cartes (genre Magic the Gathering) et jeu de plateau.

\subsubsection{L'éternel retour}
\label{sec:orge73ea48}

Au travers des jeux PbtA, le JDR vit donc, au sens strict du terme, une "régression" : là où Gary Gygax avait sorti le jeu de rôles du monde du jeu de plateau de stratégie, PbtA nous ramène vers le jeu de plateau avec les mêmes buts :
\begin{itemize}
\item Créer un jeu optimisé pour un sous-genre, voire un sous-sous-genre ;
\item Créer des archétypes de personnages taillés pour le sous-genre et aux actions limitées pour obéir aux règles du sous-genre (dans l'esprit d'une "série") ;
\item Contraindre les joueurs et le MJ à évoluer dans un narratif qui contraint les PJ dans une sandbox, dans un univers de possibles limités, où les moves ont été optimisés pour créer mécaniquement une dynamique de jeu organique.
\end{itemize}

Une des conséquences est le côté parfois extrêmement brutal ou expéditif du jeu, cela de manière très surprenante pour les pratiquants des anciens JDR. Là où les anciens JDR proposaient des approches plus progressives (difficultés pour les joueurs à mesurer les risques d'une action, mécanique d'aggravation des situations suite à échec, etc.), tout peut basculer très vite dans le monde PbtA.

En fait, ce genre de jeux est tout à fait pertinent et il est normal qu'il séduise un certain public. Mais il est aussi normal que d'autres joueurs ne l'apprécient pas, car ce n'est plus tout à fait le même \emph{type} de jeu.

\subsection{Sommes-nous vraiment tous des game designers ?}
\label{sec:orgc3ea97c}

Avec l'arrivée de plate-formes comme \href{https://www.kickstarter.com}{Kickstarter} ou \href{https://itch.io}{itch}, beaucoup de game designers se sont révélés, offrant une énorme diversité au JDR. Pour autant, la multiplication de l'offre fait apparaître des jeux \emph{dispensables} qui font se poser la question : est-ce qu'il y a autant de bons game designers sur le marché ?

Je n'en suis pas certain. Pour PbtA par exemple, les règles me semblent affreusement complexes et touffues, pleines de "trous dans la raquette" et nécessitant un investissement important pour tous les anciens MJ. Et pour quoi ? Pour faire du JDR narratif ? Mais on peut en faire facilement avec BRS ou même avec D\&D, dans des univers où les PJ sont moins caricaturaux et ont plus de possibilités de faire des choses et de s'adapter à la situation.

Pour ce qui est des scénarios ouverts, il faut se souvenir que bon nombre de scénarios anciens étaient très fouillés et très ouverts. Ils décrivaient l'univers de jeu, les PNJ, leurs motivations, le timing des événements et les PJ devaient s'insérer (voire bousculer) ces événements. Pour faire cela, il faut un système ouvert qui laisse la place belle aux inventions des joueurs, inventions qui ne manquaient pas d'arriver, souvent à la surprise du MJ.

\subsection{La manne des JDR PDF à l'heure de l'impression à la demande}
\label{sec:orga05333d}

Est-ce que le fait de lancer des systèmes complexes comme PbtA et de pouvoir lever de l'argent facilement ne rend pas plus facile l'arrivée sur le marché de produits immatures, pris dans l'engrenage financier de l'industrie du JDR ?

De plus, est-ce que les réseaux sociaux ne permettent pas de lancer des jeux dont les coûts de production sont très faibles et les revenus potentiellement très importants ? Vendre des PDF plus quelques impressions à la demande est une facilité qui collabore à mettre sur le marché des jeux parfois inaboutis.

Même chez les plus grands, le phénomène est réel. Comparons les suppléments GURPS 3e avec ceux de la 4e et vous verrez. Le digital a diminué la qualité, globalement, en poussant à la présence de suppléments, à l'exploitation du filon des suppléments sur une période courte, pendant la phase durant laquelle un jeu est à la mode.

\subsection{Le règne de D\&D 5e\ldots{} et de l'OSR}
\label{sec:org5edc7ac}

Alors, oui, dans tout cette offre pléthorique, D\&D 5e règne en maître et, semble-t-il, dans la durée. D\&D, c'est un peu l'anti-jeu moderne. Même si son système de jeu a gagné en cohérence et a pris certains éléments des nouvelles tendances, D\&D est encore D\&D.

La mode OSR (Old-School Revival) pourrait être vu comme un genre de réaction à toutes ces innovations. J'ai lu dans des forums que les OSR-guys cherchaient une façon de jouer moins complexe, mais je ne suis pas d'accord. Ils veulent du "crunchy" de la grande époque, le sommet que nous aurons du mal à dépasser : AD\&D 1e ! Des règles énormes et pleines de cas particuliers, des tables à tiroirs comme Gary Gygax les aimait, des tas et des tas d'informations de toutes sortes un peu en vrac, une certaine inventivité pour les PJ, les monstres et les pièges, un sommet du genre.

\subsection{Système de jeu idéal}
\label{sec:orgad72ad1}

Un équilibre entre :
\begin{itemize}
\item Possibilités de faire des jets de dés sous contraintes,
\item Simplicité et logique globale du système,
\item Adaptation à l'univers.
\end{itemize}

Par exemple, pour les charactéristiques, il est important qu'elles soient intuitives pour le MJ. Là dessus, D\&D et BRS sont au dessus du lot.

\subsubsection{Numéro 1 - Basic Role Playing system - BRP}
\label{sec:org756ee3d}

Le système Basic RP (\href{https://www.chaosium.com/brp-system-reference-document/}{SRD ici}), ou BRP, est un système très adaptable, logique et sans déformation de probabilités (contrairement au \href{https://github.com/orey/jdr/tree/master/D6-System}{système D6}). Il est particulièrement bien adapté aux univers fantasy, historiques et contemporains. Je ne connais pas d'implémentation du BRP en monde SF, en tous cas pas chez Chaosium.

\subsubsection{Numéro 2 - Maléfices}
\label{sec:orgb7c3419}

\begin{itemize}
\item Un système de jeu Steampunk très adapté à l'univers.
\item Tarot très utile dans le jeu.
\item Un système un peu oublié.
\end{itemize}

\subsubsection{Numéro 3 - Donjons et Dragons}
\label{sec:org544e6a9}

D\&D possède un bon système de jeu qui a fait ses preuves dans une multitude de versions. Son système est simple et basé sur le paradigme suivant : \texttt{D20 + modificateurs >= Classe de difficulté} (par exemple, dépendant plus ou moins directement de la classe d'armure). Ce système a l'avantage de ne pas tordre les probabilités (contrairement au \href{https://github.com/orey/jdr/tree/master/D6-System}{système D6}).

Voir \href{https://github.com/orey/jdr/tree/master/DandD}{la page dédiée}.

\subsection{Articles sur les probabilités dans le JDR}
\label{sec:orga410314}

\begin{itemize}
\item Une analyse des problèmes de probabilités du système D6 : voir \href{https://github.com/orey/jdr/tree/master/D6-System}{le folder D6-system}
\item Une analyse des probabilités de l'étrange système de jeu de IronSworn : voir \href{https://github.com/orey/jdr/tree/master/IronSworn}{le folder IronSworn}
\end{itemize}

\subsection{Quelques commentaires sur quelques jeux}
\label{sec:org165de9b}
\subsubsection{Méga (1984) - Méga 2 (1986)}
\label{sec:org5bb1001}
Il faut que je parle de ce jeu qui est mon premier JDR (enfin, j'avais eu D\&D avant mais je n'avais pas réussi à l'utiliser).

Un supplément fantastique publié en CC sur le \href{https://www.messagers-galactiques.com}{site de Méga IV}.

\subsubsection{Hurlements (1989)}
\label{sec:orgdbc5256}

Hurlements (1989) a quelque peu défrayé la chronique en proposant un jeu très narratif à la belle époque des jeux d'Heroic Fantasy, notamment AD\&D bien entendu. Pour autant, ce jeu ne m'a jamais convaincu, en raison, non de la pauvreté de son système de jeu, mais au niveau de la \textbf{pauvreté de sa vision du Moyen-Age}.

En effet, le jeu est centré sur la lycanthropie, mais il est très pauvre à bien des égards.

Tout d'abord, il propose une vision obscurantiste du Moyen-Age :
\begin{itemize}
\item Comme toute bonne vision caricaturale française actuelle, la religion y est caricaturée et n'est pas comprise comme un élément structurant de la société ;
\item Il n'y a aucun mot sur la chevalerie et les passions qui y sont associées, et qui sont dans la littérature française du Moyen-Age.
\end{itemize}

D'un point de vue de l'univers magique, l'univers de Hurlements est incroyablement pauvre :
\begin{itemize}
\item Ainsi, on n'y trouve pas de magie ni de sorcellerie, alors que ces éléments sont au coeur de l'univers mental de cette période ;
\item Les pouvoirs de la religion n'y sont pas évoqués.
\end{itemize}

Ainsi il aurait été intelligent de considérer qu'une certaine partie du Clergé était au courant des manifestations lycanthropiques et magiques, et que sans doute tous ne le voyaient pas forcément d'un mauvais oeil. Il aurait été intéressant par exemple de situer des abbayes comme des ponts entre les lycanthropes et certains religieux.

Concernant la magie et la sorcellerie, cette dernière était au centre de la société médiévale, tout comme l'était la chevalerie.

La vie dans la caravane est une mise en scène en mode "sandbox" qui est pourtant intéressante, mais beaucoup trop schématique (les PJ contre le reste du monde).

Je passerai sur la prétention de l'écriture qui est souvent un peu soûlante quoique très française.

Au final, ce jeu est un vrai grand raté, malgré la grosse campagne de soutien de Casus Belli et de Dragon Radieux à l'époque. Hurlements aurait pu devenir le \emph{Pendragon} français (jeu beaucoup plus mûr dans tous ses aspects) et il a sombré - assez justement - dans l'oubli.

\subsubsection{Fudge}
\label{sec:orgb88439a}

Deux éléments sont vraiment originaux :
\begin{itemize}
\item Le premier est l'usage de mots pour décrire les niveaux des caractéristiques et des skills. Fudge est sans doute un des premiers jeux à avoir fait cela (même si en fait, la mécanique de jeux reste sous-jacente et basé sur des nombres).
\item Le second est la gestion des échelles, qui est une vraie originalité de Fudge. Il est, en effet, possible de faire lutter des PJ et PNJ appartenant à des échelles différents. Une innovation très intéressante car, sur ce point, la mécanique est bien aboutie.
\end{itemize}

Voir l'\href{https://github.com/orey/jdr/tree/master/Fudge-fr}{article détaillé}.

\subsubsection{Tunnels \& Trolls}
\label{sec:org065de3a}

Le combat de groupe (mêlée) est un vrai combat de groupe :
\begin{itemize}
\item Les attaques de tous les joueurs sont cumulées,
\item Les attaques de tous les monstres le sont aussi,
\item On fait la différence (contest) pour calculer les dégâts à répartir sur la partie concernée.
\end{itemize}

Malin et efficace.

\subsubsection{Bloodlust}
\label{sec:org51c0b25}

\begin{enumerate}
\item Mécanisme de combat
\label{sec:orge1315e5}

Une seule table pour attaquant vs défenseur. En abscisse et en ordonnée :
\begin{itemize}
\item Attaque brutale
\item Attaque normale
\item Attaque rapide
\item Parade
\item Esquive
\end{itemize}

Dans le combat, chacun est tour à tour attaquant et défenseur. Fluide et efficace.

En bref, le combat est comme un double "contest" avec des modificateurs. C'est assez malin.

\item Réussites et échecs critiques
\label{sec:org1984b83}

Bloodlust est un système à pourcentage. En cas de réussite, si l'unité de la valeur du jet est 0, on est dans un cas de réussite critique. Pareil pour les échecs critiques avec une valeur de l'unité de 1 sur le jet de pourcentage raté.
\end{enumerate}
\end{document}
